% 作者题证的编辑转录副本，不是独立下载的上游原件。
% 来源：https://github.com/bbadzioch/topology_lecture_notes/blob/master/MTH_727_notes/15_weak_homotopy_type.tex

\begin{definition}[CW approximation]
A CW approximation of a space $X$ is a CW complex $Y$ together with a weak
homotopy equivalence $f:Y\to X$. More generally, a CW approximation of a pair
$(X,A)$ is a relative CW complex $(Y,A)$ and a weak equivalence $f:Y\to X$
such that $f|_A=\id_A$.
\end{definition}
\begin{theorem}[CW approximation; original label CW APPROX THM]
Any pair $(X,A)$ has a CW approximation. Any two CW approximations of this pair
are homotopy equivalent relative to $A$.
\end{theorem}
\begin{proof}
Assume first that $X$ is path connected. We construct relative CW complexes
$(Y^{(n)},A)$ and compatible maps $f^{(n)}:Y^{(n)}\to X$ so that
$Y^{(n)}$ is obtained from $Y^{(n-1)}$ by attaching $n$-cells, the maps restrict
to the identity on $A$, and
\[
 f^{(n)}_*:\pi_i(Y^{(n)})\longrightarrow\pi_i(X)
\]
is an isomorphism for $i<n$ and an epimorphism for $i=n$.
The union of the maps will give the required CW approximation.

Let $\{A_i\}_{i\in I}$ be the path components of $A$, choose $a_i\in A_i$,
and fix $x_0\in X$. Add a new $0$-cell $e^0$ to $A$. For each $i$, attach a
$1$-cell joining $e^0$ to $a_i$. For each class
$[\tau:(S^1,s_0)\to(X,x_0)]\in\pi_1(X,x_0)$, attach a circle $S^1_\tau$
by identifying $s_0$ with $e^0$. Call the resulting space $Y^{(1)}$.
Since $X$ is path connected, choose a path $\omega_i$ from $x_0$ to $a_i$.
Define $f^{(1)}$ to be the identity on $A$, to send $e^0$ to $x_0$, to map
the joining edge using $\omega_i$, and to map $S^1_\tau$ using $\tau$.
It induces a bijection on $\pi_0$ and a surjection on $\pi_1$.

Suppose now that the construction has been performed through dimension $n$.
For every element
\[
 [\omega:(S^n,s_0)\to(Y^{(n)},e^0)]\in\ker f^{(n)}_*,
\]
attach an $(n+1)$-cell using $\omega$ as attaching map. Denote the result by
$\overline Y^{(n+1)}$. Since $[f^{(n)}\omega]=0$, the composite on the
boundary extends over $D^{n+1}$, giving an extension
$\overline f^{(n+1)}:\overline Y^{(n+1)}\to X$.
Next, for each $[\tau]\in\pi_{n+1}(X,x_0)$ attach a sphere
$S^{n+1}_\tau$ at $e^0$, and extend the map by $\tau$ on that sphere.
This gives $Y^{(n+1)}$ and $f^{(n+1)}$.

With $i:Y^{(n)}\hookrightarrow Y^{(n+1)}$ the inclusion, we have
\[
\begin{tikzcd}[column sep=large]
\pi_n(Y^{(n)},e^0) \arrow[r,"i_*"] \arrow[dr,"f^{(n)}_*"'] &
\pi_n(Y^{(n+1)},e^0) \arrow[d,"f^{(n+1)}_*"]\\
&\pi_n(X,x_0).
\end{tikzcd}
\]
Surjectivity of $f^{(n)}_*$ gives surjectivity of $f^{(n+1)}_*$ in dimension
$n$, and the cells attached along its kernel make this map injective.
The lower homotopy groups are unchanged. The newly attached spheres make
$f^{(n+1)}_*$ surjective in dimension $n+1$. Thus the inductive requirements
hold. Taking $Y=\bigcup_nY^{(n)}$ and $f=\bigcup_nf^{(n)}$ gives a weak
equivalence $Y\to X$ and a relative CW structure on $(Y,A)$.

If $X$ is not path connected, let $X_i$ be its path components and apply the
construction to each pair $(X_i,A\cap X_i)$. A CW approximation of $(X,A)$ is
obtained from
\[
 A\sqcup\coprod_iY_i
\]
by identifying each point of $A\cap X_i\subset Y_i$ with the corresponding
point of $A$.

For uniqueness, suppose $f_i:(Y_i,A)\to(X,A)$ are CW approximations,
$i=1,2$. The relative homotopy lifting property for a weak equivalence and
a relative CW complex (original Corollary HELP COR) gives a map
$g:Y_1\to Y_2$ restricting to the identity on $A$ and satisfying
$f_2g\simeq f_1$ relative to $A$. Similarly, there is $h:Y_2\to Y_1$
with $f_1h\simeq f_2$ relative to $A$. Consequently there is a homotopy
$\varphi:Y_1\times I\to X$ from $f_1$ to $f_1hg$, relative to $A$.

Set
\[
 Z_1=Y_1\times\{0,1\}\cup A\times I\subset Y_1\times I.
\]
The pair $(Y_1\times I,Z_1)$ is a relative CW complex. Define
$\psi:Z_1\to Y_1$ by
\[
 \psi(y,t)=\begin{cases}y,&t<1,\\hg(y),&t=1.\end{cases}
\]
These prescriptions agree on $A\times\{1\}$ because $hg|_A=\id_A$.
We have the commutative diagram
\[
\begin{tikzcd}[column sep=large]
Z_1 \arrow[r,"\psi"] \arrow[d,hook] &Y_1\arrow[d,"f_1"]\\
Y_1\times I \arrow[r,"\varphi"'] \arrow[ur,dashed,"\overline\varphi"] & X.
\end{tikzcd}
\]
Applying the same lifting property gives $\overline\varphi:Y_1\times I\to Y_1$
extending $\psi$. It is a homotopy $\id_{Y_1}\simeq hg$ relative to $A$.
The analogous construction gives $\id_{Y_2}\simeq gh$ relative to $A$.
\end{proof}
